\section{Alternative Views}

A counter-position to our thesis holds that scaling---larger models, broader pretraining corpora, and more sophisticated world models---can eventually render runtime observation unnecessary. According to this view, L4 states are not fundamentally different from L1--L3; they are merely high-dimensional, fine-grained historical facts that a sufficiently powerful prior can approximate. If the prior captures enough of the world's dynamics, the argument goes, prediction from the prior may suffice for action, eliminating the need for real-time observation.
% 与我们论点相对的一个立场认为，规模扩展——更大的模型、更广泛的预训练语料和更精密的世界模型——最终可以使运行时观测变得不必要。根据这一观点，L4状态与L1--L3在根本上并无不同；它们仅仅是高维的、细粒度的历史事实，一个足够强大的先验可以近似它们。如果先验捕捉了足够多的世界动态，论证认为，从先验出发的预测可能就足以指导行动，从而消除了实时观测的需要。

We reject this view on three grounds. First, L4 states are \emph{indexical} by definition: they depend on \emph{when} and \emph{where} the agent is. A prior learned from past data cannot, by construction, guarantee facts about the current state that came into existence after the training cutoff. No amount of scaling can compress information that does not yet exist at training time. Second, the exponential growth of possible L4 configurations means that any finite prior will be sparse in the space of actual situated states. The agent will inevitably encounter configurations not represented in its training distribution. Third, even a perfect prior is still a distribution over states, not a certificate of the actual state. Grounding requires contact with the environment, not just probabilistic inference from memory.
% 我们基于三点理由拒绝这一观点。第一，L4状态按其定义是\emph{指示性}的：它们取决于智能体所在的\emph{何时}与\emph{何地}。从过去数据中学习的先验在构造上无法保证关于训练截止日期之后才出现的当前状态的事实。无论规模多大，都无法压缩训练时尚不存在的信息。第二，可能的L4配置呈指数增长，意味着任何有限先验在实际情境状态空间中都是稀疏的。智能体将不可避免地遇到其训练分布中未表示的配置。第三，即使是完美的先验，也仍然是状态上的分布，而非实际状态的证明。落地需要与环境接触，而不仅仅是从记忆中进行概率推理。

An alternative position might propose that embodied interaction with the environment can be folded into the LLM itself through online fine-tuning or continual learning, eliminating the need for a separate Harness architecture.
% 另一种替代立场可能提出，通过与环境的具身交互可以通过在线微调或持续学习融入LLM本身，从而消除对独立框架架构的需求。

We acknowledge that end-to-end online adaptation is an active research direction. However, our framework does not claim that the Harness is the \emph{only} possible architecture; it claims that for \emph{current} text-trained, frozen-parameter LLMs, a runtime Harness is structurally necessary. For future architectures that learn primarily through online interaction, the boundary between \(\mathcal{K}\) and \(b_t\) may shift or disappear. This is a limitation we acknowledge in Section~9.1. Our contribution is a diagnostic framework for \emph{current} systems, not a universal claim about all possible intelligent systems.
% 我们承认端到端的在线适应是一个活跃的研究方向。然而，我们的框架并不主张框架是\emph{唯一}可能的架构；它主张对于\emph{当前}的基于文本训练的、参数冻结的LLM，运行时框架在结构上是必要的。对于主要通过在线交互学习的未来架构，\(\mathcal{K}\)与\(b_t\)之间的边界可能移动或消失。这是我们在第9.1节中承认的一个局限。我们的贡献是面向\emph{当前}系统的诊断框架，而非关于所有可能智能系统的普遍主张。